%!TEX TS-program = pdflatex
%!encoding = UTF-8 Unicode
%%=============================================================================
%%=== ISEA-TEX-Vorlage ===
%%===   Verzeichnisse  ===
%%===      jmu 05.07.2012     ===
%%=============================================================================

%% 	Hier werden die Eintraege für die einzelnen Verzeichnisse definiert %%

% Zu beachten ist der type Parameter, er entscheidet in welchem Verzeichnis der Begriff erscheint

%% Umbenennung der Verzeichnisse können in Listsof.tex vorgenommen werden

%%=== Verzeichnis 1 (z.B. Abkürzungsverzeichnis) ===
\newacronym{reluk}{Geschaltete Reluktanzmaschine}{geschaltete Reluktanzmaschine}
\newacronym{rwth}{RWTH}{{Rheinisch-Westfälische Technische Hochschule Aachen}}
\newacronym{isea}{ISEA}{{Institut für Stromrichtertechnik und elektrische Antrieb}}
\newacronym{pgs}{PGS}{Institute for Power Generation and Storage Systems}
\newacronym{nmc}{NMC}{{Lithium-Nickel-Mangan-Cobalt-Oxid}}
\newacronym{lfp}{LFP}{{Lithium-Eisen-Phosphat}}
\newacronym{nco}{nco}{{Lithium-Nickel-Cobalt-Oxid}}
\newacronym{lto}{lto}{{Lithium-Titanat (Li4Ti5O12)}}
\newacronym{lipf}{$\mbox{LiPF}_6$}{{Lithiumhexafluorophosphat}}

	
%%=== Verzeichnis 2 (z.B. Symbolverzeichnis) ===


\newglossaryentry{Lambda}{type=symbolslist, name=$\Lambda$, description={Ionische Leitfähigkeit}, sort=LeitI}
	
%%=== Verzeichnis 3 (z.B. Fremdwortverzeichnis) ===

%%\newglossaryentry{Portmonaie}{name=Portmonaie}{{frz. Wort für Geldbörse}, sort=Portmonaie, type=fremdwort}


